% 09-jeu-seau.tex — Réalisation 3/4 : le jeu du seau d'eau (≈ 5 pages, rédigé).
% Implémentation : src/mechanics/loop-wagon/game.js. Les chiffres cités sont
% mesurés sur ce code (outils/donnees.mjs, tests unitaires game.test.js).

\chapter{Le jeu du seau d'eau : une corde qui plie et de l'eau en particules}
\label{chap:jeu}

\section{Principe du jeu et lien avec le looping}

Le mode jeu de l'exercice du looping remplace le wagon par une expérience que
tout le monde a vue, ou faite : faire tourner un seau d'eau à bout de bras sans
se mouiller. Une barre légère tourne librement autour d'un pivot fixe ; une corde
souple est attachée à son extrémité et porte un seau ouvert contenant dix gouttes
d'eau. Au départ, tout pend au repos ; une main balance le seau, puis le lance
(section~\ref{sec:jeu-main}) de sorte qu'il passe au point bas avec la vitesse
\(v_0\). Le défi est réussi si le seau fait deux tours sans
perdre une seule goutte et sans casser la corde ; sinon il est raté. Chaque défi
impose le lieu (Lune, Mars, Vénus, Terre, donc \(g\)) et toutes les grandeurs
sauf une, que l'étudiant règle : la vitesse \(v_0\), la longueur de corde \(L\)
ou la masse \(m\) du seau rempli.

Tant que la corde est tendue, barre, corde et seau restent alignés et forment
un pendule de rayon \(R = L_b + L + d\), où \(L_b = \SI{0,50}{m}\) est la longueur
de la barre et \(d = \SI{0,19}{m}\) la distance de l'anse au centre de masse du
seau. La conservation de l'énergie donne \(v^2 = v_0^2 - 2 g R (1 - \cos(\theta))\),
\(\theta\) étant l'angle depuis le point bas, et la deuxième loi de Newton selon
la direction radiale donne la tension
\begin{equation}
  T = m \left(\frac{v^2}{R} + g \cos(\theta)\right)
    = m \left(\frac{v_0^2}{R} + g \, (3 \cos(\theta) - 2)\right).
  \label{eq:jeu-tension}
\end{equation}
C'est exactement la réaction normale du wagon, \(N = m (v^2/r + g \cos(\varphi))\),
avec \(v_0^2 = 2 g h\) (tableau~\ref{tab:jeu-analogie}). Pourquoi l'eau reste-t-elle
dans le seau ? Seuls le poids \(\vect{P}\) et la tension \(\vect{T}\) agissent sur
le seau, donc \(\vect{a} = \vect{g} + \vect{T}/m\). Une goutte libre a pour
accélération \(\vect{g}\) ; relativement au seau, elle accélère donc de
\(-\vect{T}/m\), c'est-à-dire vers le fond du seau tant que \(T > 0\). L'eau
reste dans le seau si et seulement si la corde reste tendue, et la condition la
plus sévère est au sommet : \(T_s = m (v_0^2/R - 5 g) \geq 0\), soit
\(v_0^2 \geq 5 g R\). La corde résiste si la tension maximale, au point bas, ne
dépasse pas la tension de rupture : \(m (v_0^2/R + g) \leq T_{\max}\).

\begin{table}[ht]
  \centering\small
  \begin{tabular}{@{}lll@{}}
    \toprule
    & Wagon sur le looping & Seau au bout de la corde \\
    \midrule
    Force de contact & réaction normale \(N\) & tension \(T\) \\
    Énergie disponible & \(v_{\text{bas}}^2 = 2 g h\) & \(v_0^2\) \\
    Contact perdu si & \(N < 0\) : décollage du rail & \(T = 0\) : corde molle, l'eau tombe \\
    Condition au sommet & \(h \geq 5R/2\) & \(v_0^2 \geq 5 g R\) \\
    Limite supplémentaire & --- & \(m (v_0^2/R + g) \leq T_{\max}\) \\
    \bottomrule
  \end{tabular}
  \caption{Le jeu reprend la physique de l'exercice : la condition de réussite est
    la relation du cours, pas un réglage à l'aveugle.}
  \label{tab:jeu-analogie}
\end{table}

\paragraph{Pourquoi une barre libre et non un moteur ?} Une barre entraînée à
vitesse angulaire \(\omega\) constante donne un comportement chaotique. Dans le
référentiel tournant, le seau pendu à une corde de longueur \(\ell\) fixée à la
distance \(a\) du pivot est un pendule centrifuge (\(\Omega = \omega \sqrt{a/\ell}\))
forcé une fois par tour par la pesanteur ; l'amplitude forcée, environ
\(g / (\omega^2 |a - \ell|)\), est de l'ordre de \SI{1}{rad} près du seuil. Une barre
légère sur un pivot sans frottement ne transmet au contraire aucun couple : la
corde tendue s'aligne sur elle et la relation~\eqref{eq:jeu-tension} vaut.

\section{Un modèle discret : une chaîne de points}

Le système est représenté par une chaîne de treize points matériels
(figure~\ref{fig:jeu-schema}). Le pivot \(P_0\) est fixe. L'extrémité de la barre
\(P_1\) porte la masse effective \(m_b/3\), qui reproduit le moment d'inertie
\(I = m_b L_b^2/3\) d'une barre homogène de masse \(m_b = \SI{0,3}{kg}\) ; on y applique
le poids \(m_b g/2\), pour que le couple du poids, \(m_b g \, (L_b/2) \sin(\alpha)\), soit
exact. La corde
est découpée en dix maillons de longueur \(\ell = L/10\) séparés par les nœuds
\(P_2\) à \(P_{11}\), chacun de masse \(\mu \ell\) avec une masse linéique
\(\mu = \SI{0,15}{kg/m}\) ; \(P_{11}\) est l'anse du seau et \(P_{12}\) son centre
de masse, de masse \(m\). Deux liaisons sont rigides dans les deux sens : la barre
\(|P_1 - P_0| = L_b\) et le seau \(|P_{12} - P_{11}| = d\), qui fait du seau un
solide articulé à son anse. Les dix maillons de corde sont \emph{unilatéraux} :
\(|P_{i+1} - P_i| \leq \ell\). Un maillon peut raccourcir, jamais s'allonger, et
aucun angle n'est imposé entre deux maillons. C'est ce qui rend la corde
pliable : tendue, elle est droite ; détendue, elle se courbe et ne pousse pas le
seau.

\begin{figure}[ht]
  \centering
  \begin{tikzpicture}[x=1cm, y=1cm, >={Stealth[length=2mm]},
      noeud/.style={circle, fill=corde, inner sep=1.3pt},
      etiquette/.style={font=\scriptsize, align=center}]
    \fill[metal!60] (-0.25,-0.35) -- (0.25,-0.35) -- (0,0) -- cycle;
    \draw[line width=3pt, metal] (0,0) -- (2.2,0);
    \fill[black] (0,0) circle (3pt);
    \node[etiquette, above] at (0,0.1) {\(P_0\)\\pivot fixe};
    \node[etiquette, below] at (1.1,-0.15) {barre rigide \(L_b\)};
    \foreach \k in {0,...,10} {
      \coordinate (c\k) at ({2.2 + 0.5*\k}, {-0.55*sin(18*\k)});
    }
    \draw[corde, line width=1.4pt] (c0) \foreach \k in {1,...,10} { -- (c\k) };
    \foreach \k in {1,...,9} { \node[noeud] at (c\k) {}; }
    \fill[metal] (c0) circle (2.4pt);
    \node[etiquette, above] at (c0) {\(P_1\)\\\(m_b/3\)};
    \node[etiquette, below] at (c5) {\(P_2 \ldots P_{10}\) : nœuds de corde \(\mu \ell\)\\maillons \(|P_{i+1} - P_i| \leq \ell\)};
    \draw[metal, line width=1.2pt, fill=metal!15] (7.6,-0.65) -- (8.85,-0.58) -- (8.85,0.58) -- (7.6,0.65) -- cycle;
    \draw[thin] (7.6,0.65) .. controls (7.0,0.45) and (7.0,-0.45) .. (7.6,-0.65);
    \foreach \x/\y in {8.68/-0.42, 8.68/-0.14, 8.68/0.14, 8.68/0.42, 8.44/-0.28, 8.44/0, 8.44/0.28, 8.2/-0.14, 8.2/0.14, 7.96/0} {
      \fill[eau] (\x,\y) circle (0.12);
    }
    \fill[black] (7.2,0) circle (2pt);
    \node[etiquette, above left] at (7.2,0.05) {\(P_{11}\)\\anse};
    \fill[poids] (8.4,0) circle (2pt);
    \node[etiquette, right] at (8.95,0.25) {\(P_{12}\) : centre\\de masse, \(m\)};
    \draw[->, accent, thick] (7.2,0) -- (7.2,-1.1) node[below, etiquette] {\(\vect{n}\)};
    \draw[->, accent, thick] (8.4,-0.95) -- (9.3,-0.95) node[right, etiquette] {\(\vect{e}\)};
    \draw[|-|, thin] (7.2,1.0) -- node[above, etiquette] {\(u_{\text{bord}}\)} (7.6,1.0);
    \draw[|-|, thin] (7.2,1.45) -- node[above, etiquette] {\(u_{\text{fond}} = \SI{26}{cm}\)} (8.85,1.45);
  \end{tikzpicture}
  \caption{Modèle discret : pivot fixe, barre (masse effective \(m_b/3\)), dix
    maillons de corde unilatéraux, seau rigide articulé à son anse. À droite, le
    repère du seau : axe \(\vect{e}\) de l'anse vers le fond, normale \(\vect{n}\),
    bord à \(u_{\text{bord}} = \SI{6}{cm}\) et fond à \(u_{\text{fond}} = \SI{26}{cm}\)
    de l'anse ; les dix gouttes reposent sur le fond.}
  \label{fig:jeu-schema}
\end{figure}

\section{L'élan : une main qui porte le seau}
\label{sec:jeu-main}

Un élan réaliste ne peut pas commencer par une corde tendue et un seau déjà
lancé : la corde ne doit jamais être figée. Au départ, tout pend au repos ; une
main tient l'anse \(P_{11}\) et la déplace sur un cercle de rayon
\((L_b + L)(1 - s)\), avec un mou \(s = \SI{3}{\percent}\) : la corde, un peu détendue,
se courbe et entraîne la barre libre (figure~\ref{fig:jeu-captures}). La main fait
d'abord un petit va-et-vient (\(+\ang{12}\) puis \(-\ang{45}\), à peu près au rythme
propre du pendule, \(\tau = \sqrt{R/g}\)), puis un lancer par en dessous jusqu'à
l'angle \(\theta_r\) où elle lâche le seau : \(+\ang{90}\) en général, \(+\ang{150}\) pour
les lancers rapides (\(v_0^2 \geq 6 g R\)), dont le trajet plus long adoucit les
accélérations. Pour que l'eau ne déborde pas pendant ce transport, la main incline
le seau selon la pesanteur apparente \(\vect{g} - \vect{a}\) de l'anse, comme on
porte un seau plein.

Le lancer suit un polynôme de degré~6 en temps. Il part du repos sans
accélération et arrive en \(\theta_r\) avec la vitesse
\(v_r = \sqrt{v_0^2 - 2 g R (1 - \cos(\theta_r))}\), de sorte que le seau passe
ensuite au point bas à \(v_0\) par conservation de l'énergie, ainsi qu'avec
l'accélération et la dérivée de l'accélération du mouvement libre en ce point.
\(\vect{g} - \vect{a}\) y est donc radial : le seau est aligné sur la corde et la
passation se fait sans à-coup ; le mou n'est plus que de \SI{0,2}{\percent}. La
dérivée troisième compte : l'inclinaison du seau suit \(\vect{a}\), donc sa vitesse
de rotation suit la dérivée de \(\vect{a}\), et un saut de celle-ci au lâcher
changerait d'environ \SI{0,7}{m/s} la vitesse du centre du seau. Comme le seau
tourne avec son axe, son centre ne va pas exactement à \(\omega r\) : la vitesse
angulaire finale de la main est corrigée par quelques itérations sur la vitesse
du centre, calculée à partir de la seule trajectoire de la main. Au lâcher, le
centre part ainsi à \SI{0,2}{\percent} près de \(v_r\). Pendant l'élan, la corde ne porte que la barre
et la rupture n'est testée qu'après le lâcher. Deux pistes plus
\enquote{automatiques} ont été écartées : une barre menée selon un horaire fixé
(la vitesse au lâcher varie de façon chaotique avec la vitesse de la main) et
une main qui pilote la barre par rétroaction sur l'énergie du seau (à-coups de
tension, corde qui fouette).

\section{Dynamique basée sur les positions}

Modéliser une corde inextensible par des ressorts très raides imposerait un pas
de temps minuscule. La dynamique basée sur les positions \parencite{muller2007pbd}
contourne le problème : on prédit d'abord le mouvement sans liaisons, puis on
\emph{projette} les positions prédites sur les liaisons, et l'on déduit enfin les
vitesses des déplacements (algorithme~\ref{alg:jeu-pas}). Sa forme étendue XPBD
\parencite{macklin2016xpbd} associe à chaque liaison un multiplicateur de Lagrange
\(\lambda\), dont on tire directement la force : la tension d'un maillon vaut
\(T = -\lambda/\Delta t^2\), et l'on retient celle du maillon le plus chargé (côté barre
pour une corde lourde). Nous utilisons l'approche à petits pas
\parencite{macklin2019smallsteps} : \num{2400} sous-pas par seconde après le lâcher
(\num{1200} pendant l'élan, où la corde est molle) et une seule projection par
sous-pas, plus précis que quelques gros pas itérés.

\begin{algorithm}[ht]
  \caption{Un sous-pas de la simulation du seau (\(\Delta t = 1/2400\) s).}
  \label{alg:jeu-pas}
  \begin{algorithmic}[1]
    \For{chaque point libre \(i\) (nœuds, seau, gouttes)}
      \State \(\vect{x}^{\,\text{préc}}_i \gets \vect{x}_i\) ;\quad
             \(\vect{v}_i \gets \vect{v}_i + \vect{g} \, \Delta t\) ;\quad
             \(\vect{x}_i \gets \vect{x}_i + \vect{v}_i \, \Delta t\)
    \EndFor
    \State projeter toutes les liaisons de la chaîne en une fois (section~\ref{sec:jeu-direct}) ; \(T \gets -\lambda_{\text{corde}}/\Delta t^2\)
    \State ramener chaque goutte dans le seau, séparer les gouttes (section~\ref{sec:jeu-eau})
    \For{chaque point libre \(i\)}
      \State \(\vect{v}_i \gets (\vect{x}_i - \vect{x}^{\,\text{préc}}_i)/\Delta t\)
    \EndFor
    \State filtrer \(T\) sur \SI{20}{ms} ; corde cassée si \(T > T_{\max}\) ; goutte perdue si elle franchit le bord
  \end{algorithmic}
\end{algorithm}

\section{Faire tenir une corde pliable : la résolution directe}
\label{sec:jeu-direct}

La méthode habituelle projette les liaisons une par une (Gauss--Seidel). Elle
échoue ici à cause du rapport des masses : un nœud de corde pèse quelques
grammes, le seau jusqu'à plusieurs kilogrammes. Chaque projection déplace
presque uniquement le nœud léger, et la traction du seau ne remonte la chaîne
que d'un maillon par itération. Avec une seule passe par sous-pas, une corde de
\SI{30}{cm} portant un seau de \SI{5}{kg} s'allonge de plus de \SI{80}{cm}
(figure~\ref{fig:jeu-verification}a), ce qui fausse complètement la tension.

Nous résolvons donc toutes les liaisons simultanément. Linéarisées, les \(n\)
liaisons \(C_i = |\vect{x}_{i+1} - \vect{x}_i| - \ell_i\) donnent un système
\(A \, \Delta\boldsymbol{\lambda} = -\boldsymbol{C}\). Comme chaque point n'est
partagé que par deux liaisons voisines, la matrice est tridiagonale :
\begin{equation}
  A_{i,i} = w_i + w_{i+1}, \qquad
  A_{i,i+1} = A_{i+1,i} = -w_{i+1} \, (\vect{n}_i \cdot \vect{n}_{i+1}),
\end{equation}
où \(w_i = 1/m_i\) est l'inverse de la masse (nul pour le pivot) et \(\vect{n}_i\)
le vecteur unitaire du maillon \(i\). L'algorithme de Thomas
\parencite{press2007} la résout en \(O(n)\) opérations, puis chaque point est
corrigé : \(\vect{x}_i \mathrel{-}= w_i \vect{n}_i \Delta\lambda_i\) et
\(\vect{x}_{i+1} \mathrel{+}= w_{i+1} \vect{n}_i \Delta\lambda_i\). Cette idée de
solveur direct pour les chaînes est celle de \textcite{deul2018direct}. Deux
linéarisations par sous-pas suffisent : l'allongement résiduel reste inférieur à
\SI{e-5}{mm}.

Reste le caractère unilatéral. Une corde ne peut que tirer : pour chaque maillon,
\(C_i \leq 0\), \(\lambda_i \leq 0\) et \(\lambda_i C_i = 0\) (soit il est tendu, soit
il ne transmet rien). C'est un problème de complémentarité linéaire, résolu par
une méthode d'ensemble actif \parencite{nocedal2006} : tous les maillons partent
actifs ; on libère celui qui pousserait le plus (\(\Delta\lambda_i > 0\)) ; on reprend
un maillon inactif que la correction allongerait ; on recommence jusqu'à ce que
plus rien ne change. Un dernier piège est apparu en mettant au point le code :
avec une corde trop légère, les ondes transversales \(c = \sqrt{T/\mu}\) parcourent
un maillon plus vite qu'un sous-pas, et la corde vibre de façon parasite. La
masse linéique réaliste de \SI{0,15}{kg/m} règle le problème.

\section{L'eau : dix particules}
\label{sec:jeu-eau}

Chaque goutte est un disque de rayon \(r = \SI{2,2}{cm}\) soumis à son seul poids.
À chaque sous-pas, on calcule le repère du seau à partir de l'anse \(P_{11}\) et du
centre de masse \(P_{12}\) : axe \(\vect{e}\) et normale \(\vect{n}\)
(figure~\ref{fig:jeu-schema}). Les coordonnées locales d'une goutte sont
\(u = (\vect{x} - \vect{x}_{11}) \cdot \vect{e}\) et
\(w = (\vect{x} - \vect{x}_{11}) \cdot \vect{n}\). La goutte est ramenée dans le seau
par projection : \(u \leq u_{\text{fond}} - r\) pour le fond et
\(|w| \leq \SI{10}{cm} - r\) pour les parois. Deux gouttes ne peuvent pas se
chevaucher (\(|\vect{x}_i - \vect{x}_j| \geq 2 r\), correction partagée par moitié),
ce qui suffit pour qu'elles s'empilent au fond. Sans frottement, dix particules
glisseraient bien plus qu'une eau réelle au moindre basculement : leur vitesse
relative à la paroi est amortie (taux \SI{8}{s^{-1}}, une viscosité effective),
ce qui ne change pas le critère \(T \geq 0\). Une goutte dont le centre franchit
le plan du bord (\(u < u_{\text{bord}}\)) est perdue : elle continue sa chute libre
jusqu'au sol. Le couplage est à sens unique : les gouttes ne poussent pas le seau,
leur masse est comprise dans \(m\). Cette simplification assumée ne change pas la
condition de réussite, puisque celle-ci ne dépend que de la tension.

La figure~\ref{fig:jeu-captures} montre un élan complet puis un essai raté sur la
Lune. Quand \(v_0\) est trop faible, la corde se détend avant le sommet ; seau et
eau tombent ensemble, presque en chute libre, et l'eau n'est éjectée que lorsque
la corde se retend brutalement. Ce scénario n'a pas été programmé : il découle des
seules liaisons.

\begin{figure}[ht]
  \centering
  \newcommand{\vignette}[2]{%
    \begin{subfigure}[b]{0.325\linewidth}
      \includegraphics[width=\linewidth]{figures/captures/#1}
      \caption*{#2}
    \end{subfigure}}
  \vignette{jeu-elan-1.png}{élan : petit va-et-vient}\hfill
  \vignette{jeu-elan-3.png}{fin du retour à \(-\ang{45}\)}\hfill
  \vignette{jeu-elan-4.png}{lancer : la corde molle plie}\\[3pt]
  \vignette{jeu-elan-6.png}{seau lâché, corde tendue}\hfill
  \vignette{jeu-rate-2.png}{\(v_0\) trop faible : corde molle}\hfill
  \vignette{jeu-rate-4.png}{l'eau est perdue}
  \caption{Captures du jeu (Lune, \(L = \SI{0,89}{m}\), \(m = \SI{4,7}{kg}\),
    \(T_{\max} = \SI{78}{N}\)). Les quatre premières vignettes : \(v_0 = \SI{4,18}{m/s}\),
    réussi. Les deux dernières : \(v_0 = \SI{3,04}{m/s} = 0{,}85\,\sqrt{5 g R}\), raté.
    La main (beige) tient l'anse pendant l'élan ; poids \(P\) en rouge, tension \(T\)
    en violet après le lâcher.}
  \label{fig:jeu-captures}
\end{figure}

\section{Validation et jouabilité}

La simulation doit donner raison à la formule du cours, sinon le jeu
contredirait l'exercice. Pour 24 configurations (4 lieux, \(L\) = 0,3 ; 0,7 ;
\SI{1,2}{m}, \(m\) = 1 et \SI{5}{kg}), le seuil de réussite simulé, trouvé par
dichotomie, vaut \(v_0^\star = k \sqrt{5 g R}\) avec \(k\) compris entre \num{0,962}
et \num{1,002}, élan compris. Le seuil est un peu sous la valeur théorique parce que
la barre et la corde, lancées avec le seau, apportent leur propre énergie
cinétique, que la formule à corde sans masse ignore. Au point bas, la tension
simulée dépasse \(m (v_0^2/R + g)\) d'environ \SI{5}{\percent}
(figure~\ref{fig:jeu-verification}b) : c'est la charge de la corde elle-même,
\(\mu L (\omega^2 r_{\text{milieu}} + g)\). La projection des positions dissipe en outre
un peu d'énergie (environ \SI{5}{\percent} en deux tours), un effet du premier ordre
en \(\Delta t\).

\begin{figure}[ht]
  \centering
  \begin{subfigure}[b]{0.49\linewidth}
    \begin{tikzpicture}
      \begin{axis}[width=\linewidth, height=4.6cm, xlabel={\(t\) (s)}, ylabel={allongement (mm)},
        xmin=0, xmax=0.8, grid=major, grid style={gray!20}, legend style={font=\scriptsize},
        legend pos=north east, tick label style={font=\scriptsize}, label style={font=\small}]
        \addplot[poids, thick] table[x=t, y=gauss_seidel, col sep=comma] {figures/donnees/allongement_corde.csv};
        \addplot[accent, very thick] table[x=t, y=directe, col sep=comma] {figures/donnees/allongement_corde.csv};
        \legend{{Gauss--Seidel (une passe)}, résolution directe}
      \end{axis}
    \end{tikzpicture}
    \caption{Allongement total d'une corde de \SI{30}{cm} (seau de \SI{5}{kg}).}
  \end{subfigure}\hfill
  \begin{subfigure}[b]{0.49\linewidth}
    \begin{tikzpicture}
      \begin{axis}[width=\linewidth, height=4.6cm, xlabel={\(t\) (s)}, ylabel={\(T\) (N)},
        xmin=0, grid=major, grid style={gray!20}, legend style={font=\scriptsize},
        legend pos=north east, tick label style={font=\scriptsize}, label style={font=\small}]
        \addplot[color=tension, thick] table[x=t, y=simulee, col sep=comma] {figures/donnees/tension_seau.csv};
        \addplot[black, dashed] table[x=t, y=ideale, col sep=comma] {figures/donnees/tension_seau.csv};
        \legend{simulée, \(m (v_0^2/R + g (3 \cos(\theta) - 2))\)}
      \end{axis}
    \end{tikzpicture}
    \caption{Tension après le lâcher (Terre, \(L = \SI{0,8}{m}\), \(v_0 = 1{,}15 \sqrt{5 g R}\)).}
  \end{subfigure}
  \caption{Vérifications numériques, calculées avec le code de l'application.}
  \label{fig:jeu-verification}
\end{figure}

Chaque défi doit être \emph{prouvablement} soluble. La génération est
constructive : on tire le lieu, \(L\) et \(m\), puis une vitesse de référence
\(v_0^2 = (5{,}8 \text{ à } 7{,}5) \, g R\), et l'on fixe
\(T_{\max} = m (v_0^2/R + g) \times (1{,}25 \text{ à } 1{,}6)\). La fenêtre gagnante
annoncée est la fenêtre idéale réduite de \SI{6}{\percent} de chaque côté sur \(v_0\) ;
sa borne haute tient aussi compte de la charge propre de la corde. Elle est
convertie si nécessaire en fenêtre sur \(L\) ou sur \(m\) (pour la masse, il faut
remplir le seau à au moins \SI{90}{\percent} de la masse maximale). Le vérificateur
\fichier{isChallengeFeasible} ne se contente pas d'une inégalité : il \emph{rejoue
la simulation complète}, élan compris, en cinq points régulièrement répartis dans
la fenêtre, et le défi n'est accepté que si les cinq essais réussissent. Un essai
complet prend de 70 à \SI{150}{ms} sur Terre et la génération d'un défi environ
\SI{0,25}{s}. Les tests unitaires (dont un tirage de Monte-Carlo de 500 défis)
vérifient ces propriétés.

\begin{pointcle}
\textbf{Limites.} Le modèle est plan, sans frottement au pivot ni résistance de
l'air, avec une petite dissipation numérique. La main est idéalisée : sa
trajectoire est imposée et l'inclinaison du seau néglige l'accélération angulaire
propre du seau, ce qui ne compte que pour des lancers bien plus violents que ceux
des défis. L'eau n'a pas de pression, sa viscosité est une simple traînée, et
elle ne réagit pas sur le seau. Un fluide
plus réaliste (SPH ou fluides basés sur les positions,
\cite{monaghan1992sph, macklin2013pbf}) est possible, mais il n'apporterait rien à
la question posée à l'étudiant : à quelle vitesse l'eau sort-elle du seau ?
\end{pointcle}
